% Source: ACM_Templates.pdf, PDF pp. 69-71.
% Review and validation: docs/validation.md.
\section{最小生成树：Kruskal 与 Prim}
\label{algo:mst}
最小生成树连接无向图的全部顶点，使用 $n-1$ 条边且总权最小。若图不连通，则不存在生成树。两种算法都基于切分性质：跨过某个切分的最小权边，可以加入某棵最小生成树。边权可以为负，重边允许，自环不会被选中。

\subsection*{Kruskal 原理与使用教程}
按边权从小到大枚举，用并查集判断两个端点是否属于不同连通块；只有连接不同块时才选择该边。选中 $n-1$ 条才得到生成树，否则当前结果只是最小生成森林。时间 $O(m\log(m+2))$，空间 $O(n+m)$。

先复制 \InlineCode{code/data-structures/dsu.cpp}。\InlineCode{add(u,v,w)} 每条无向边只加一次，返回从 $0$ 开始的边编号；\InlineCode{run()} 返回是否连通，\InlineCode{cost} 为总权，\InlineCode{chosen} 为选中边编号。树权和及累加中间量须在 64 位范围内。
\begin{lstlisting}[style=mtycpp,numbers=none]
Kruskal g;
g.init(3);
g.add(1,2,2),g.add(2,3,3),g.add(1,3,8);
bool ok=g.run();
if(ok){long long ans=g.cost;}
\end{lstlisting}

\needspace{7\baselineskip}
\subsection*{模板代码}
\CodeFile{code/graphs/kruskal.cpp}

\subsection*{Prim 原理与使用教程}
从一个点出发，每次选择连接“已加入点集”和“未加入点集”的最小权边。本实现把候选边放入小根堆，弹出已加入的顶点时跳过。全部点被加入即成功，时间保守为 $O((n+m)\log(m+2))$，空间 $O(n+m)$。

\InlineCode{add(u,v,w)} 内部会加两个方向，不要重复输入反向边。\InlineCode{run(s)} 从 $s$ 出发，默认 $1$；返回真后读取 \InlineCode{cost}。要求 $n\geq1$；单个顶点的树权为 $0$。
\begin{lstlisting}[style=mtycpp,numbers=none]
Prim g;
g.init(3);
g.add(1,2,-2),g.add(2,3,3),g.add(1,3,8);
bool ok=g.run();
if(ok){long long ans=g.cost;}
\end{lstlisting}

\needspace{7\baselineskip}
\subsection*{模板代码}
\CodeFile{code/graphs/prim.cpp}

稠密图还可用邻接矩阵实现 $O(n^2)$ 的朴素 Prim：每次线性寻找最小连接代价，再更新其余点。Prim 维护的是接入树的最小边权，而非源点到该点的路径长度，不要与 Dijkstra 混用。
